\documentclass[11pt]{article}
\usepackage[T1]{fontenc}
\usepackage[margin=1in]{geometry}
\usepackage{amsmath,amssymb,microtype}
\usepackage[colorlinks=true,urlcolor=blue,citecolor=blue]{hyperref}
\title{The later C$^*$-uniqueness answer to arXiv:2107.05806}
\author{Literature-status note for \texttt{fa\_banach\_001}}
\date{11 August 2026}

\begin{document}
\maketitle

\noindent\textbf{Status.}
\texttt{literature\_already\_answered} (Question 9.2: necessary and
sufficient conditions for uniqueness of a C$^*$-norm).

\section*{The exact source question}

Mori's paper introduces R$^*$-algebras---algebraic pre-C$^*$-algebras, or
equivalently regular $^*$-subalgebras of bounded operators---and proves that
every R$^*$-algebra has a unique C$^*$-norm.  In its Questions section it asks:
\begin{quote}
Does the study of R$^*$-algebras give a better understanding of the
uniqueness of C$^*$-norm?  In particular, can we obtain a nice
necessary/sufficient condition for a $^*$-algebra to have a unique C$^*$-norm?
\end{quote}
This is Question 9.2 of the published numbering \cite{Mori}.

\section*{The later affirmative answer}

An and Gao explicitly reproduce this as their Question 2 and state that
Section 4 answers it affirmatively \cite{AG}.  Their results give not merely
one criterion but several equivalent descriptions for a pre-C$^*$-algebra
$\mathcal A$:
\begin{itemize}
  \item Theorem 4.1 characterizes C$^*$-uniqueness through the
  $^*$-representations of $\mathcal A$.
  \item Theorem 4.4 gives an equivalent formulation through ideals of the
  enveloping C$^*$-algebra $C^*(\mathcal A)$.
  \item Theorem 4.6 reformulates it in the primitive ideal space of
  $C^*(\mathcal A)$, and Theorem 5.6 gives a weak-containment criterion for
  representations.
\end{itemize}
There is also a concrete spectral test.  Theorem 4.7 says that a unital
pre-C$^*$-algebra $\mathcal A$ has a unique C$^*$-norm if and only if, for
every pair of injective $^*$-representations
\[
 \pi_i:\mathcal A\longrightarrow B(H_i)\quad(i=1,2)
\]
and every self-adjoint $a\in\mathcal A$, the operator $\pi_1(a)$ is
invertible exactly when $\pi_2(a)$ is invertible.  Thus the invertibility of
self-adjoint elements is representation-independent.  Theorem 4.10 provides
the corresponding R$^*$-algebra formulation, recovering uniqueness of the
C$^*$-norm for every R$^*$-algebra.

\section*{Conclusion and scope}

This is an exact literature resolution of the concrete existence question:
the later paper names Mori's question, answers it affirmatively, and supplies
multiple necessary-and-sufficient conditions.  It does not settle the other
independent questions in Mori's Section 9 (for example the Kadison finite
generation problem, tensor products of hypothetical non-ultramatricial
R$^*$-algebras, or the noncommutative Kalton--Roberts problem).  The result
classification here is therefore deliberately restricted to Question 9.2.

\begin{thebibliography}{9}
\bibitem{Mori}
M.~Mori,
\emph{On regular $^*$-algebras of bounded linear operators: A new approach
towards a theory of noncommutative Boolean algebras},
arXiv:2107.05806,
\url{https://arxiv.org/abs/2107.05806}.

\bibitem{AG}
G.~An and M.~Gao,
\emph{R$^*$-algebras and C$^*$-uniqueness of pre-C$^*$-algebras},
Journal of Mathematical Analysis and Applications 526 (2023), 127341,
\url{https://doi.org/10.1016/j.jmaa.2023.127341}.
\end{thebibliography}

\end{document}
